% Copyright 2021 Politecnico di Milano, Italy. NC-BY

\documentclass{Configuration_Files/PoliMi3i_thesis}

%------------------------------------------------------------------------------
%	REQUIRED PACKAGES AND  CONFIGURATIONS
%------------------------------------------------------------------------------

% CONFIGURATIONS
\usepackage{parskip} % For paragraph layout
\usepackage{setspace} % For using single or double spacing
\usepackage{emptypage} % To insert empty pages
\usepackage{multicol} % To write in multiple columns (executive summary)
\setlength\columnsep{15pt} % Column separation in executive summary
\setlength\parindent{0pt} % Indentation
\raggedbottom  

% PACKAGES FOR TITLES
\usepackage{titlesec}
% \titlespacing{\section}{left spacing}{before spacing}{after spacing}
\titlespacing{\section}{0pt}{3.3ex}{2ex}
\titlespacing{\subsection}{0pt}{3.3ex}{1.65ex}
\titlespacing{\subsubsection}{0pt}{3.3ex}{1ex}
\usepackage{color}

% PACKAGES FOR LANGUAGE AND FONT
\usepackage[english]{babel} % The document is in English  
\usepackage[utf8]{inputenc} % UTF8 encoding
\usepackage[T1]{fontenc} % Font encoding
\usepackage[11pt]{moresize} % Big fonts

% PACKAGES FOR IMAGES
\usepackage{graphicx}
\usepackage{transparent} % Enables transparent images
\usepackage{eso-pic} % For the background picture on the title page
\usepackage{subfig} % Numbered and caption subfigures using \subfloat.
\usepackage{tikz} % A package for high-quality hand-made figures.
\usetikzlibrary{}
\graphicspath{{./Images/}} % Directory of the images
\usepackage{caption} % Coloured captions
\usepackage{xcolor} % Coloured captions
\usepackage{amsthm,thmtools,xcolor} % Coloured "Theorem"
\usepackage{float}

% STANDARD MATH PACKAGES
\usepackage{amsmath}
\usepackage{amsthm}
\usepackage{amssymb}
\usepackage{amsfonts}
\usepackage{bm}
\usepackage[overload]{empheq} % For braced-style systems of equations.
\usepackage{fix-cm} % To override original LaTeX restrictions on sizes

% PACKAGES FOR TABLES
\usepackage{tabularx}
\usepackage{longtable} % Tables that can span several pages
\usepackage{colortbl}

% PACKAGES FOR ALGORITHMS (PSEUDO-CODE)
\usepackage{algorithm}
\usepackage{algorithmic}

% PACKAGES FOR REFERENCES & BIBLIOGRAPHY
\usepackage[colorlinks=true,linkcolor=black,anchorcolor=black,citecolor=black,filecolor=black,menucolor=black,runcolor=black,urlcolor=black]{hyperref} % Adds clickable links at references
\usepackage{cleveref}
\usepackage[square, numbers, sort&compress]{natbib} % Square brackets, citing references with numbers, citations sorted by appearance in the text and compressed
\bibliographystyle{abbrvnat} % You may use a different style adapted to your field

% OTHER PACKAGES
\usepackage{pdfpages} % To include a pdf file
\usepackage{afterpage}
\usepackage{lipsum} % DUMMY PACKAGE
\usepackage{fancyhdr} % For the headers
\fancyhf{}

% Input of configuration file. Do not change config.tex file unless you really know what you are doing. 
% Define blue color typical of polimi
\definecolor{bluepoli}{cmyk}{0.4,0.1,0,0.4}

% Custom theorem environments
\declaretheoremstyle[
  headfont=\color{bluepoli}\normalfont\bfseries,
  bodyfont=\color{black}\normalfont\itshape,
]{colored}

% Set-up caption colors
\captionsetup[figure]{labelfont={color=bluepoli}} % Set colour of the captions
\captionsetup[table]{labelfont={color=bluepoli}} % Set colour of the captions
\captionsetup[algorithm]{labelfont={color=bluepoli}} % Set colour of the captions

\theoremstyle{colored}
\newtheorem{theorem}{Theorem}[chapter]
\newtheorem{proposition}{Proposition}[chapter]

% Enhances the features of the standard "table" and "tabular" environments.
\newcommand\T{\rule{0pt}{2.6ex}}
\newcommand\B{\rule[-1.2ex]{0pt}{0pt}}

% Pseudo-code algorithm descriptions.
\newcounter{algsubstate}
\renewcommand{\thealgsubstate}{\alph{algsubstate}}
\newenvironment{algsubstates}
  {\setcounter{algsubstate}{0}%
   \renewcommand{\STATE}{%
     \stepcounter{algsubstate}%
     \Statex {\small\thealgsubstate:}\space}}
  {}

% New font size
\newcommand\numfontsize{\@setfontsize\Huge{200}{60}}

% Title format: chapter
\titleformat{\chapter}[hang]{
\fontsize{30}{20}\selectfont\bfseries\filright}{\textcolor{bluepoli} \thechapter\hsp\hspace{2mm}\textcolor{bluepoli}{|   }\hsp}{0pt}{\huge\bfseries \textcolor{bluepoli}
}

% Title format: section
\titleformat{\section}
{\color{bluepoli}\normalfont\Large\bfseries}
{\color{bluepoli}\thesection.}{1em}{}

% Title format: subsection
\titleformat{\subsection}
{\color{bluepoli}\normalfont\large\bfseries}
{\color{bluepoli}\thesubsection.}{1em}{}

% Title format: subsubsection
\titleformat{\subsubsection}
{\color{bluepoli}\normalfont\large\bfseries}
{\color{bluepoli}\thesubsubsection.}{1em}{}

% Shortening for setting no horizontal-spacing
\newcommand{\hsp}{\hspace{0pt}}

\makeatletter
% Renewcommand: cleardoublepage including the background pic
%\renewcommand*\cleardoublepage{%
%  \clearpage\if@twoside\ifodd\c@page\else
%  \null
%  \AddToShipoutPicture*{\BackgroundPic}
%  \thispagestyle{empty}%
%  \newpage
%  \if@twocolumn\hbox{}\newpage\fi\fi\fi}
\makeatother

%For correctly numbering algorithms
\numberwithin{algorithm}{chapter}

\usepackage{etoolbox}
\makeatletter
% Se nel wrapper della classe c'è un \vspace{-1cm}, sostituiscilo:
\patchcmd{\polimichapter}{\vspace{-1cm}}{\vspace{0pt}}{}{}
\patchcmd{\polimichapterstar}{\vspace{-1cm}}{\vspace{0pt}}{}{}
\makeatother

\usepackage{titlesec}
% \titlespacing*{\chapter}{<rientro>}{<spazio PRIMA>}{<spazio DOPO>}
\titlespacing*{\chapter}{0pt}{10pt}{8pt} % es.: niente spazio sopra, 20pt sotto

% --- Niente salto pagina all'inizio capitolo, ma mantiene numerazione/TOC ---
\usepackage{etoolbox}
\makeatletter
% rimuove il clearpage/cleardoublepage dentro \chapter (stile 'book')
\patchcmd{\chapter}{\if@openright\cleardoublepage\else\clearpage\fi}{}{}{}
% se la classe ha salvato il \chapter in \latexchapter e lo richiama via wrapper, patchiamo anche quello (se esiste)
\patchcmd{\latexchapter}{\if@openright\cleardoublepage\else\clearpage\fi}{}{}{}
\makeatother


%----------------------------------------------------------------------------
%	NEW COMMANDS DEFINED
%----------------------------------------------------------------------------

% EXAMPLES OF NEW COMMANDS
\newcommand{\bea}{\begin{eqnarray}} % Shortcut for equation arrays
\newcommand{\eea}{\end{eqnarray}}
\newcommand{\e}[1]{\times 10^{#1}}  % Powers of 10 notation

%----------------------------------------------------------------------------
%	ADD YOUR PACKAGES (be careful of package interaction)
%----------------------------------------------------------------------------

%----------------------------------------------------------------------------
%	ADD YOUR DEFINITIONS AND COMMANDS (be careful of existing commands)
%----------------------------------------------------------------------------

%----------------------------------------------------------------------------
%	BEGIN OF YOUR DOCUMENT
%----------------------------------------------------------------------------

\begin{document}

\fancypagestyle{plain}{%
\fancyhf{} % Clear all header and footer fields
\fancyhead[RO,RE]{\thepage} %RO=right odd, RE=right even
\renewcommand{\headrulewidth}{0pt}
\renewcommand{\footrulewidth}{0pt}}

%----------------------------------------------------------------------------
%	TITLE PAGE
%----------------------------------------------------------------------------

\pagestyle{empty} % No page numbers
\frontmatter % Use roman page numbering style (i, ii, iii, iv...) for the preamble pages

\puttitle{
	title={Convolutional Neural Networks: Theory, Architectures and Applications}, % Title of the thesis
	name=Giorgio Monaco, % Author Name and Surname
	course=Numerical Analysis for Machine Learning, % 
	ID  = 277930,  % Student ID number (numero di matricola)
	advisor= Prof. Edie Miglio, % Supervisor name
	coadvisor=, % Co-Supervisor name, remove this line if there is none
	academicyear={2024-25},  % Academic Year
} % These info will be put into your Title page 

%----------------------------------------------------------------------------
%	LIST OF CONTENTS/FIGURES/TABLES/SYMBOLS
%----------------------------------------------------------------------------

% TABLE OF CONTENTS
\thispagestyle{empty}
\tableofcontents % Table of contents 
\thispagestyle{empty}
\cleardoublepage

%-------------------------------------------------------------------------
%	THESIS MAIN TEXT
%-------------------------------------------------------------------------
% In the main text of your thesis you can write the chapters in two different ways:
%
%(1) As presented in this template you can write:
%    \chapter{Title of the chapter}
%    *body of the chapter*
%
%(2) You can write your chapter in a separated .tex file and then include it in the main file with the following command:
%    \chapter{Title of the chapter}
%    \input{chapter_file.tex}
%
% Especially for long thesis, we recommend you the second option.

\addtocontents{toc}{\vspace{2em}} % Add a gap in the Contents, for aesthetics
\mainmatter % Begin numeric (1,2,3...) page numbering

% --------------------------------------------------------------------------
% NUMBERED CHAPTERS % Regular chapters following
% --------------------------------------------------------------------------

\chapter*{Introduction}

Deep learning has transformed the way machines perceive and interpret data, with Convolutional Neural Networks (CNNs) emerging as one of the most powerful and widely adopted architectures, particularly in computer vision. Unlike traditional fully connected networks, which struggle with high-dimensional inputs such as images, CNNs exploit spatial hierarchies in data to efficiently extract relevant features, making them the backbone of most modern visual recognition systems.
\\
The origins of CNNs trace back to the 1990s, when Yann LeCun and colleagues introduced \textit{LeNet-5}, a pioneering architecture for handwritten digit recognition \cite{lecun1998gradient}. Although LeNet demonstrated the potential of convolutional layers, limited computational resources at the time restricted its widespread adoption. The true breakthrough arrived in 2012 with \textit{AlexNet} \cite{krizhevsky2012imagenet}, which leveraged GPU acceleration, ReLU activations, and dropout regularization to win the ImageNet Large Scale Visual Recognition Challenge, reducing the top-5 error rate from 26\% to 15\% and igniting the modern deep learning era.
\\
Since then, CNNs have evolved rapidly. \textit{VGGNet} \cite{simonyan2014very} showcased the benefits of deep architectures with uniform $3 \times 3$ convolutions, while \textit{ResNet} \cite{he2015deep} introduced residual connections to enable the training of extremely deep networks. More recently, \textit{EfficientNet} \cite{tan2019efficientnet} optimized model scaling strategies, achieving state-of-the-art performance with fewer computational resources.
\\
Today, CNNs underpin a wide range of applications, from image classification and object detection to medical imaging, natural language processing, and autonomous robotics. Despite their success, CNNs still face challenges such as high computational costs, limited interpretability, and vulnerability to adversarial inputs \cite{li2021survey}. Nevertheless, their impact on artificial intelligence remains transformative, and their role is expected to endure even as hybrid approaches with emerging architectures, such as Transformers, gain traction.
\\
This paper provides an introduction to CNNs by exploring their mathematical foundations, architectural components, training procedures, landmark models, and real-world applications. By the end, readers should gain a comprehensive understanding of how CNNs work, why they are effective, and where the field is heading.

\chapter{Theoretical and Mathematical Foundations}
\label{ch:th_math_foundations}%

Understanding the theoretical foundations of Convolutional Neural Networks (CNNs) requires revisiting the basic concepts of artificial neurons, fully connected layers, and the convolution operation itself.


\section{Artificial Neuron and Fully-Connected Layers}
\label{sec:an_fc_layers_recap}%

An artificial neuron receives inputs \(x_1, x_2,...,x_n\), multiplies them by learnable weights \(w_1, w_2,...,w_n\), sums them together, adds a bias term \(b\), and applies a non-linear activation function \(\sigma(\cdot)\):

\begin{align}
    z = \sum_{i=1}^{n} w_ix_i + b \label{eq:ann_1} \\
    y = \sigma(x) \label{eq:ann_2}
\end{align}

In a fully connected layer (or dense layer), every neuron in a given layer is connected to all neurons in the previous layer. While this structure is powerful, it does not scale well to high-dimensional inputs such as images.
\\
For example, a single RGB image of size \(224 \times 224\) has: \(224\times224\approx150.000\) input features.
\\
A fully connected network would require millions of parameters, making it inefficient and prone to overfitting: CNNs overcome this limitation through the convolution operation.

\section{The Convolution Operation}
\label{sec:convolution_op}%

\subsection{Mathematical Definition}
\label{subsec:math_definition}

A convolution is a mathematical operation that combines two functions to produce a third function expressing how one shape modifies the other.

\begin{itemize}
    \item Continuous convolution (1D):

    \begin{align}
    (f*g)(t)=\int_{-\infty}^{+\infty}f(\tau)g(t-\tau)d\tau
    \end{align}

    \item Discrete convolution (1D):

    \begin{align}
    (f*g)[n]=\sum_{m=-\infty}^{+\infty}f[m]g[n-m]
    \end{align}
    
\end{itemize}

In CNNs, we usually work with discrete 2D convolutions for image data. For an image \(I\) and a kernel (or filter) 
\(K\), the convolution at location (\(i,j\)) is:

\begin{align}
    S(i,j) = (I*K)(i,j)=\sum_m\sum_nI(i+m,j+n)\cdot K(m,n)
\end{align}


\subsection{1D vs 2D Convolutions}

\begin{itemize}
    \item 1D convolutions are often used in time-series analysis or signal processing.
    Example: detecting patterns in ECG signals or sound waves.

    \item 2D convolutions are used in image processing, where filters (kernels) slide across height and width to extract spatial features.
\end{itemize}

\subsection{Numerical Example}
Consider a small \(3\times3\) grayscale image \(I\) and a \(2\times2\) kernel \(K\):

\begin{align}
    I = \begin{bmatrix} 1 & 2 & 3 \\ 4 & 5 & 6 \\ 7 & 8 & 9 \end{bmatrix}, K = \begin{bmatrix} 1 &  0 \\ 0 & -1 \end{bmatrix}
\end{align}

The convolution at the top-left position is:

\begin{align}
    S(0,0)=1\cdot1+2\cdot0+4\cdot0+5\cdot(-1)=1-5=-4
\end{align}

By sliding the kernel across the image, we obtain the feature map \(S\).
\\
This process highlights patterns such as edges or textures depending on the kernel values.


\subsection{Convolution vs Cross-Correlation}

In mathematics, convolution flips the kernel before applying it. However, in practice, most deep learning frameworks (e.g., TensorFlow, PyTorch) implement cross-correlation, defined as:

\begin{align}
    S(i,j)=\sum_m\sum_n I(i+m,j+n)\cdot K(m,n)
\end{align}

The difference is subtle but important:
\begin{itemize}
    \item Convolution: kernel is reversed.
    \item Cross-correlation: kernel is used as-is.
\end{itemize}

In CNNs, we learn the kernel weights, so the flip does not affect the final result.

\section{Data Representation in CNNs}

CNNs operate directly on tensors, which are multi-dimensional arrays of data.

\subsection{Kernels, Feature Maps, and Receptive Fields}
A kernel (or filter) is a small learnable matrix, that extracts specific features such as edges, textures, or shapes. The feature map is the result of convolving an image with a kernel, highlighting regions where the feature is present. The receptive field refers to the region of the input image that influences a particular neuron in the feature map.
\\
As we go deeper into the network, neurons have larger receptive fields and capture higher-level features, from simple edges in the first layers to complex shapes and objects in deeper layers.
\\
For example, imagine sliding a \(3\times3\) kernel across a \(5\times5\) image:
\\
At each step, the kernel multiplies corresponding pixel values and sums them up. The resulting feature map shows where patterns detected by the kernel occur in the image.
\\
Multiple kernels can be stacked in a single convolutional layer to extract different types of features simultaneously.

\chapter{Basic CNN Architecture}
\label{ch:basic_cnn_arch}%

A \textbf{Convolutional Neural Network (CNN)} is composed of a sequence of specialized layers designed to efficiently process high-dimensional inputs such as images.
\\
Unlike fully connected networks, CNNs exploit the spatial structure of data to extract hierarchical features while keeping the number of learnable parameters relatively low \cite{goodfellow2016deep}.
\\
In this section, we explore the main building blocks of a CNN: \textbf{convolutional layers}, \textbf{pooling layers}, \textbf{activation functions}, \textbf{batch normalization}, and \textbf{fully connected layers}.

\section{Convolutional Layers}

The \textbf{convolutional layer} is the core component of a CNN. It applies a set of learnable \textbf{kernels} (or filters) to the input image or feature map, extracting patterns such as \textbf{edges}, \textbf{textures}, or \textbf{shapes} \cite{goodfellow2016deep}.
\\
Each kernel slides across the input with a given stride and computes a dot product between the kernel weights and the local region of the input, producing a feature map.
This process allows the network to detect local structures while sharing parameters across the entire image, dramatically reducing the total number of learnable weights compared to fully connected architectures.
\\
Key hyperparameters of a convolutional layer:

\begin{itemize}
    \item \textbf{Kernel size}: determines the size of the receptive field, e.g. $3 \times 3$ or $5 \times 5$.
    \item \textbf{Stride} : controls how far the kernel moves at each step; larger strides reduce the spatial resolution.
    \item \textbf{Padding}: determines how to handle edges.

    \begin{itemize}
        \item Valid padding: no padding, output shrinks.
        \item Same padding: zero-padding to maintain input dimensions.
    \end{itemize}
\end{itemize}

For example, a convolutional layer with 64 filters of size $3 \times 3$ on a $224 \times 224 \times 3$ input produces 64 feature maps, each detecting a different learned pattern.

\section{Pooling Layers}

Pooling layers are used to downsample the spatial dimensions of feature maps while retaining the most important information \cite{goodfellow2016deep}.
This reduces the computational cost and improves the model’s ability to generalize.
\\
The two most common types are:

\begin{itemize}
    \item \textbf{Max} pooling: selects the maximum value within a local region.
    \item \textbf{Average} pooling: computes the average value within the region.
\end{itemize}

For instance, applying a $2 \times 2$ max pooling operation with stride 2 on a $32 \times 32$ feature map reduces the output to $16 \times 16$.
\\
Pooling layers offer two main advantages:
\begin{enumerate}
    \item Translation invariance: small shifts in the input do not significantly affect the feature maps.
    \item Dimensionality reduction: fewer parameters and less computation, reducing the risk of overfitting.
\end{enumerate}

\section{Activation Functions}

Activation functions introduce \textbf{non-linearity} into the network, enabling CNNs to learn complex patterns that cannot be captured by linear transformations alone.
Without them, stacking multiple layers would result in a model equivalent to a single linear transformation.
\\
Common activation functions include:

\begin{itemize}
    \item \textbf{ReLU} (Rectified Linear Unit):
    
    \begin{align}
        \text{ReLU}(x) = \max(0, x)
    \end{align}
    
    Computationally efficient and alleviates the vanishing gradient problem.

    \item \textbf{Leaky} ReLU: allows a small, non-zero gradient for negative inputs, addressing “dying ReLU” neurons.

    \item \textbf{ELU} (Exponential Linear Unit): smooths the activation curve and sometimes improves learning speed.
    
\end{itemize}

According to Li et al. \cite{li2021survey}, ReLU and its variants are by far the most widely used activation functions in modern CNNs because of their simplicity and effectiveness.

\section{Batch Normalization}

\textbf{Batch Normalization (BatchNorm)} normalizes the activations of a layer across each mini-batch, ensuring they have zero mean and unit variance \cite{li2021survey}.
This technique has several benefits:
\begin{itemize}
    \item Faster convergence during training.
    \item Stabilized learning, reducing sensitivity to weight initialization.
    \item Regularization effect, improving generalization and reducing overfitting.
\end{itemize}

Introduced in 2015, BatchNorm is now a standard component in most modern CNNs and is typically placed after the convolution but before the activation function.

\section{Fully Connected Layers}

After several convolutional and pooling layers, the extracted features are flattened and passed to one or more fully connected (dense) layers.
These layers integrate the learned features and perform the final classification or regression task.
\\
Each neuron in a fully connected layer connects to all activations from the previous layer. The final layer often uses a softmax activation to produce a probability distribution over output classes.
\\
For example, in the ImageNet classification task, the last fully connected layer outputs a 1,000-dimensional vector, where each value corresponds to the predicted probability of a specific class \cite{goodfellow2016deep}.

\section{Summary of the Data Flow}

A typical CNN processes data as follows:
\begin{enumerate}
    \item Input Layer: raw image data.
    \item Convolutional Layers: extract local patterns such as edges and textures.
    \item Pooling Layers: reduce spatial dimensions while preserving essential information.
    \item Batch Normalization \& Activations: stabilize and accelerate learning.
    \item Fully Connected Layers: combine extracted features into final predictions.
    \item Output Layer: provides class probabilities or regression values, depending on the task.
\end{enumerate}

\chapter{Training Convolutional Neural Networks}

Training a Convolutional Neural Network (CNN) involves teaching the network to extract meaningful features from input data and map them to the desired outputs.
This process relies on two fundamental steps: the forward pass, where predictions are computed, and the backpropagation, where errors are propagated backward to update the network’s parameters \cite{goodfellow2016deep}.
\\
In this chapter, we explain how CNNs learn through:

* The forward propagation of information.
* The calculation of loss functions.
* The optimization of parameters using gradient-based methods.
* Techniques to avoid overfitting and improve generalization.

\section{Forward Pass}

During the forward pass, the input data flows sequentially through the network’s layers:
\\
1. Convolutional layers apply filters to extract hierarchical features.
2. Activation functions introduce non-linearity, enabling the network to model complex relationships.
3. Pooling layers reduce spatial dimensions while retaining critical information.
4. Fully connected layers integrate the learned features to make predictions.
5. The output layer produces final values, whether they are probabilities for classification tasks or continuous values for regression.
\\
For example, in an image classification task on ImageNet, the forward pass transforms an input image into a 1,000-dimensional probability vector, where each dimension corresponds to a possible class.

\section{Backpropagation and Weight Updates}

After the forward pass, the CNN evaluates how far its predictions are from the expected output using a loss function.
The network then uses backpropagation to update its weights and minimize this error \cite{goodfellow2016deep}.
\\
The process involves:
\\
1. Computing the gradient of the loss with respect to each parameter using the chain rule.
2. Propagating the error backward from the output layer through the network.
3. Updating weights based on these gradients using an optimization algorithm.
\\
Backpropagation is crucial in CNNs, as it allows the network to learn which filters are useful for recognizing specific patterns.
Through many iterations, the kernels gradually specialize in detecting relevant features at different levels of abstraction.

\section{Loss Functions}

Loss functions measure the discrepancy between the predicted outputs and the ground truth labels.
The choice of loss function depends on the task:
\\
\textbf{Cross-Entropy Loss}, commonly used for classification tasks:

\begin{align}
    \mathcal{L} = - \sum_{i=1}^C y_i \cdot \log(\hat{y}_i)
\end{align}

where:

\begin{itemize}
    \item $C$ is the number of classes,
    \item $y_i$ is the true label (one-hot encoded),
    \item $\hat{y}_i$ is the predicted probability for class $i$.
\end{itemize}

\textbf{Mean Squared Error (MSE)}, used for regression tasks:

\begin{align}
    \mathcal{L} = \frac{1}{N} \sum_{i=1}^N (y_i - \hat{y}_i)^2
\end{align}

where $N$ is the number of samples.
\\
Choosing the right loss function is essential for guiding the network to learn meaningful representations \cite{goodfellow2016deep}.

\section{Optimization Algorithms}

To minimize the loss, CNNs rely on gradient-based optimization methods.
Some of the most widely used include:
\begin{itemize}
    \item \textbf{Stochastic Gradient Descent (SGD)}:
        Updates weights based on small batches of data, it's simple but effective, especially with momentum to accelerate convergence.
    \item \textbf{Adam} (Adaptive Moment Estimation):
        Combines the benefits of SGD with momentum and adaptive learning rates. Popular in modern architectures due to fast convergence and robustness.
    \item \textbf{RMSProp}:
        Maintains an adaptive learning rate for each parameter, it's useful for tasks with non-stationary objectives, such as recurrent models.
\end{itemize}

According to Li et al. \cite{li2021survey}, Adam and SGD with momentum remain the most widely adopted optimizers in CNNs due to their balance between efficiency and stability.

\chapter{Landmark and Modern Architectures}

Since the introduction of the first Convolutional Neural Networks, several \textbf{landmark architectures} have shaped the evolution of deep learning. Each new generation addressed the limitations of its predecessors, enabling networks to become deeper, more efficient, and more accurate on large-scale datasets such as ImageNet.
\\
In this chapter, we analyze five key architectures that defined the development of CNNs: \textbf{LeNet-5}, \textbf{AlexNet}, \textbf{VGGNet}, \textbf{ResNet}, and \textbf{EfficientNet}.

\section{LeNet-5 (1998)}

Developed by Yann LeCun and colleagues, LeNet-5 \cite{lecun1998gradient} was one of the first CNN architectures designed for handwritten digit recognition. It achieved excellent results on the MNIST dataset, reaching nearly perfect accuracy, and was even deployed in banking systems for automated check reading.
\\
LeNet-5 introduced several concepts that remain foundational in modern CNNs, such as convolutional layers, pooling operations, and feature maps.
\\
The architecture processed $32 \times 32$ grayscale images through two convolutional layers, each followed by average pooling, and concluded with two fully connected layers for classification. Although modest in scale compared to modern networks, LeNet laid the groundwork for all subsequent developments.

\section{AlexNet (2012)}

The real breakthrough came in 2012 when Alex Krizhevsky, Ilya Sutskever, and Geoffrey Hinton introduced AlexNet \cite{krizhevsky2012imagenet}. Trained on the ImageNet dataset, AlexNet reduced the top-5 error rate from 26\% to 15\%, outperforming all competing models by a significant margin and marking the beginning of the deep learning revolution.
\\
AlexNet’s success stemmed from several innovations. It popularized the use of the ReLU activation function, which accelerated convergence compared to sigmoid and tanh. It also introduced dropout regularization to combat overfitting and leveraged GPU acceleration to train a large model efficiently. Furthermore, data augmentation techniques such as image translations and flips were used to artificially expand the training set, further improving generalization.

\section{VGGNet (2014)}

Two years later, Karen Simonyan and Andrew Zisserman proposed VGGNet \cite{simonyan2014very}, which became renowned for its simplicity and depth. Unlike AlexNet, VGGNet adopted a uniform architecture, stacking multiple $3 \times 3$ convolutional filters sequentially. This design enabled the network to reach depths of 16 or 19 layers (VGG16 and VGG19, respectively), significantly improving performance on ImageNet.
\\
However, this gain in accuracy came at a cost: VGGNet required an enormous number of parameters (approximately 138 million for VGG16) making it computationally expensive and memory intensive. Despite this drawback, VGGNet remains a popular feature extractor in transfer learning due to its well-structured representation of visual patterns.

\section{ResNet (2015)}

As networks grew deeper, researchers encountered the vanishing gradient problem, where gradients became too small to effectively update early layers during training. To address this issue, Kaiming He and colleagues introduced ResNet \cite{he2015deep}, a groundbreaking architecture that implemented residual connections.
\\
These skip connections allow layers to learn residual mappings rather than direct transformations. Mathematically, a residual block learns:
\begin{align}
    y = F(x) + x
\end{align}

This innovation enabled the successful training of extremely deep models, such as ResNet-152, without performance degradation. ResNet achieved state-of-the-art results on ImageNet, winning the ILSVRC 2015 competition and inspiring a new family of architectures based on residual learning.

\section{EfficientNet (2019)}

In 2019, Mingxing Tan and Quoc V. Le introduced EfficientNet \cite{tan2019efficientnet}, which proposed a compound scaling method to balance three dimensions of a CNN simultaneously: depth, width, and input resolution. Previous models typically scaled these dimensions independently, leading to suboptimal designs. EfficientNet’s compound scaling achieved a better trade-off between accuracy and efficiency.
\\
EfficientNet models achieved state-of-the-art performance on ImageNet with significantly fewer parameters and FLOPs compared to previous architectures. For instance, EfficientNet-B7 reached a top-1 accuracy of nearly 97\% while using fewer computational resources than much larger models.

\section{Comparative Overview}

The following table summarizes the key characteristics of these landmark architectures:

\begin{table}[H]
\centering
\begin{tabular}{lcccccl}
\hline
\textbf{Architecture} & \textbf{Year} & \textbf{Depth} & \textbf{Params} & \textbf{Dataset} & \textbf{Top-1 Acc.} & \textbf{Key Innovation} \\
\hline
LeNet-5 & 1998 & 7  & $\sim$60K  & MNIST    & $\sim$99\%  & First CNN, pooling layers \\
AlexNet & 2012 & 8  & 61M        & ImageNet & 84.7\%      & ReLU, dropout, GPU training \\
VGG16   & 2014 &16  & 138M       & ImageNet & $\sim$90\%  & Small $3\times3$ kernels, deep uniform design \\
ResNet50 & 2015 &50 & 25.6M      & ImageNet & $\sim$96\%  & Residual connections \\
EfficientNet-B7 & 2019 &66 & 66M & ImageNet & $\sim$97\%  & Compound scaling \\
\hline
\end{tabular}
\caption{Comparison of landmark CNN architectures.}
\label{tab:cnn_comparison}
\end{table}

The evolution of CNNs reflects a constant balance between accuracy, efficiency, and scalability.
LeNet-5 laid the foundations by introducing convolution and pooling.
AlexNet triggered the deep learning revolution through GPU training and ReLU activations.
VGGNet demonstrated the effectiveness of deeper yet simpler architectures, while ResNet solved the training challenges of very deep models using residual connections.
Finally, EfficientNet redefined model scaling strategies, achieving state-of-the-art performance with fewer computational resources \cite{li2021survey}.

\chapter{Practical Applications of CNNs}

Convolutional Neural Networks (CNNs) have become the standard for a wide range of tasks in artificial intelligence, especially in the field of computer vision. 
Their ability to automatically extract hierarchical features from raw data has enabled breakthroughs across domains such as image recognition, object detection, medical imaging, and autonomous systems \cite{li2021survey, zhao2024review}.
\\
In this section, we explore the most relevant applications of CNNs, highlighting how the architectures presented in Section~\ref{tab:cnn_comparison} have been successfully deployed in real-world scenarios.

\section{Image Classification}

One of the earliest and most impactful applications of CNNs has been \textbf{image classification}, where the goal is to assign a label to an entire image. 
Starting from LeNet-5's success on MNIST to AlexNet's breakthrough on the ImageNet dataset, CNNs have consistently pushed the boundaries of classification performance.
\\
Modern architectures such as VGG, ResNet, and EfficientNet achieve top-1 accuracies above 90\% on ImageNet, making CNNs the backbone of most commercial image recognition systems \cite{li2021survey}. 
These capabilities power numerous services, from automatic photo tagging on social media platforms to visual search engines and large-scale multimedia retrieval systems.

\section{Object Detection and Localization}

While classification focuses on identifying what is present in an image, \textbf{object detection} extends the task to also determining \textit{where} objects are located.
CNN-based models such as R-CNN, Faster R-CNN, YOLO, and SSD leverage feature maps to predict bounding boxes and class labels simultaneously. 
\\
Thanks to these advances, CNNs are now widely used in applications such as:

\begin{itemize}
    \item \textbf{Autonomous driving}: detecting pedestrians, vehicles, and traffic signs in real time.
    \item \textbf{Security systems}: monitoring surveillance footage for abnormal activities.
    \item \textbf{Retail analytics}: recognizing products and tracking customer behavior.
\end{itemize}

According to Zhao et al. \cite{zhao2024review}, the integration of CNNs with attention mechanisms has further improved detection accuracy, especially in challenging environments.

\section{Semantic and Instance Segmentation}

Another critical application of CNNs is \textbf{image segmentation}, which involves partitioning an image into regions corresponding to different objects or classes.
Two main approaches are commonly distinguished:

\begin{itemize}
    \item \textbf{Semantic segmentation}: assigns a class label to every pixel, grouping identical objects under the same label.
    \item \textbf{Instance segmentation}: goes a step further by differentiating between individual objects of the same class.
\end{itemize}

Architectures such as U-Net, Mask R-CNN, and DeepLab have become standard tools in this field \cite{zhao2024review}. 
Applications range from autonomous navigation to industrial quality control and satellite image analysis.

\section{Medical Imaging}

Medical image analysis has benefited enormously from the adoption of CNNs, which are now routinely used in:

\begin{itemize}
    \item Detecting tumors and lesions in MRI and CT scans.
    \item Classifying skin diseases based on dermoscopic images.
    \item Automating radiological workflows by assisting with diagnosis.
\end{itemize}

According to Li et al. \cite{li2021survey}, CNN-based systems in medical imaging have achieved human-level accuracy in several diagnostic tasks. 
These advances accelerate clinical decision-making and reduce the workload of healthcare professionals.

\section{Autonomous Driving and Robotics}

Self-driving cars and autonomous robots rely heavily on CNNs to interpret their surroundings. 
Through a combination of image classification, object detection, and segmentation, CNNs enable real-time understanding of complex environments \cite{zhao2024review}. 
\\
For instance, Tesla’s Autopilot and Waymo’s self-driving systems integrate CNNs to detect pedestrians, vehicles, and obstacles, as well as to understand lane markings and traffic lights.

\section{Beyond Computer Vision}

Although CNNs were originally developed for visual tasks, they have been successfully applied beyond computer vision:

\begin{itemize}
    \item \textbf{Natural Language Processing}: character-level CNNs for text classification and sentiment analysis.
    \item \textbf{Audio Processing}: spectrogram-based CNNs for speech recognition and music genre classification.
    \item \textbf{Multimodal Models}: combining CNNs with Transformers for tasks that integrate vision and language, such as image captioning.
\end{itemize}

\chapter{Challenges and Future Directions}

Despite their remarkable success, Convolutional Neural Networks (CNNs) still face several open challenges that limit their broader applicability. 
\\
As highlighted in recent surveys \cite{li2021survey, zhao2024review}, current research focuses on improving the efficiency, robustness, and interpretability of these models while exploring new paradigms that may complement or even surpass CNNs.

\section{Computational and Memory Constraints}

Modern CNN architectures, such as VGG and EfficientNet, often involve tens of millions of parameters and require significant computational resources during both training and inference. 
This makes their deployment difficult on mobile devices, embedded systems, and other environments where hardware is limited.  
\\
To address this limitation, researchers are developing approaches such as model compression, parameter pruning, and quantization to reduce memory usage and computational cost without sacrificing accuracy. Another promising direction is knowledge distillation, where smaller “student” models are trained to mimic the behavior of larger, more complex “teacher” networks, making CNNs more suitable for real-time applications \cite{zhao2024review}.

\section{Interpretability and Explainability}

Although CNNs achieve state-of-the-art results on many vision tasks, they are often criticized for being black-box models.
\\
Understanding why a model makes a particular prediction is crucial in domains where transparency is required, such as healthcare, finance, and autonomous driving.  
Recent work has introduced techniques like saliency maps, class activation maps (CAMs), and feature attribution methods to provide visual and quantitative explanations of network decisions. 
\\
As Li et al. \cite{li2021survey} observe, improving interpretability will be essential to fostering trust in AI systems, especially in sensitive and safety-critical contexts.

\section{Robustness and Adversarial Attacks}

Another major challenge is robustness. CNNs perform well on clean datasets but are vulnerable to adversarial attacks — small, carefully designed perturbations that can mislead the model into making incorrect predictions with high confidence.  
This vulnerability poses significant risks in applications like autonomous driving, biometric authentication, and medical diagnostics.  
\\
To mitigate this issue, researchers are exploring adversarial training, defensive distillation, and certified defenses, but achieving robust and secure CNNs remains an open problem \cite{zhao2024review}.

\section{Beyond CNNs: The Rise of Vision Transformers}

While CNNs dominate computer vision, a new class of models known as Vision Transformers (ViTs) has recently gained attention.  
Unlike CNNs, which rely on local receptive fields and convolutional operations, ViTs use self-attention mechanisms to model global dependencies across an entire image.  
\\
This shift has led to significant improvements in benchmark performance, especially when large datasets are available.  
However, ViTs often require more computational resources, which has spurred interest in hybrid architectures combining the strengths of CNNs and Transformers \cite{zhao2024review}.

\section{Future Perspectives}

Looking ahead, research is likely to focus on making CNNs more lightweight, interpretable, and robust.  
Efforts to optimize architectures for mobile and edge devices will make deep learning more accessible in real-time applications.  
\\
At the same time, developing transparent and explainable models will be key to gaining acceptance in sensitive areas such as healthcare and security.  
Finally, the integration of CNNs with emerging paradigms — including Vision Transformers and graph neural networks — suggests a future where hybrid architectures become the standard \cite{li2021survey, zhao2024review}.

\chapter{Conclusion and Outlook}

Convolutional Neural Networks have transformed machine learning by combining mathematical structure, computational efficiency and inspiration from neuroscience.  
Through convolution, they exploit local connectivity, weight sharing and translation equivariance, while pooling introduces useful invariances and reduces complexity.  
Variants of the basic convolution function extend their flexibility and modern alternatives to pooling demonstrate the adaptability of these models to evolving needs.  
Applications across one, two and three-dimensional data, highlight the generality of convolution as a tool for representation learning.  
The neuroscientific parallels further reinforce the conceptual foundation of CNNs, even though many differences remain between artificial models and biological vision.

Looking forward, convolutional networks are likely to remain central in deep learning, though their dominance is being challenged by newer architectures such as attention based models and transformers.  
Future directions may involve hybrid models that combine convolution with attention or other mechanisms, better approximations of biological processes such as foveation and feedback and broader applications beyond vision, including multimodal learning and scientific data analysis.  

In summary, CNNs illustrate how mathematical principles, engineering solutions and biological insights can converge into models that not only achieve state of the art performance but also continue to inspire new generations of machine learning research.



\chapter{Chapter one}
\label{ch:chapter_one}%
% The \label{...}% enables to remove the small indentation that is generated, always leave the % symbol.

In this chapter additional useful information are reported.

\section{Sections and subsections}
\label{sec:section_name}
Chapters are typically subdivided into sections and subsections, and, optionally,
subsubsections, paragraphs and subparagraphs.
All can have a title, but only sections and subsections are numbered.
A new section is created by the command
\begin{verbatim}
\section{Title of the section}
\end{verbatim}
The numbering can be turned off by using \verb|\section*{}|.
\\
A new subsection is created by the command
\begin{verbatim}
\subsection{Title of the subsection}
\end{verbatim}
and, similarly, the numbering can be turned off by adding an asterisk as follows 
\begin{verbatim}
\subsection*{}
\end{verbatim}

\section{Equations}
\label{sec:eqs}
This section gives some examples of writing mathematical equations in your thesis.

Maxwell's equations read:
\begin{subequations}
    \label{eq:maxwell}
    \begin{align}[left=\empheqlbrace]
    \nabla\cdot \bm{D} & = \rho, \label{eq:maxwell1} \\
    \nabla \times \bm{E} +  \frac{\partial \bm{B}}{\partial t} & = \bm{0}, \label{eq:maxwell2} \\
    \nabla\cdot \bm{B} & = 0, \label{eq:maxwell3} \\
    \nabla \times \bm{H} - \frac{\partial \bm{D}}{\partial t} &= \bm{J}. \label{eq:maxwell4}
    \end{align}
\end{subequations}

Equation~\eqref{eq:maxwell} is automatically labeled by \texttt{cleveref},
as well as Equation~\eqref{eq:maxwell1} and Equation~\eqref{eq:maxwell3}.
Thanks to the \verb|cleveref| package, there is no need to use \verb|\eqref|.
Remember that Equations have to be numbered only if they are referenced in the text.

Equations~\eqref{eq:maxwell_multilabels1}, \eqref{eq:maxwell_multilabels2}, \eqref{eq:maxwell_multilabels3}, and \eqref{eq:maxwell_multilabels4} show again Maxwell's equations without brace:
\begin{align}
    \nabla\cdot \bm{D} & = \rho, \label{eq:maxwell_multilabels1} \\
    \nabla \times \bm{E} +  \frac{\partial \bm{B}}{\partial t} &= \bm{0}, \label{eq:maxwell_multilabels2} \\
    \nabla\cdot \bm{B} & = 0, \label{eq:maxwell_multilabels3} \\
    \nabla \times \bm{H} - \frac{\partial \bm{D}}{\partial t} &= \bm{J} \label{eq:maxwell_multilabels4}.
\end{align}

Equation~\eqref{eq:maxwell_singlelabel} is the same as before,
but with just one label:
\begin{equation}
    \label{eq:maxwell_singlelabel}
    \left\{
    \begin{aligned}
    \nabla\cdot \bm{D} & = \rho, \\
    \nabla \times \bm{E} +  \frac{\partial \bm{B}}{\partial t} &= \bm{0},\\
    \nabla\cdot \bm{B} & = 0, \\
    \nabla \times \bm{H} - \frac{\partial \bm{D}}{\partial t} &= \bm{J}.
    \end{aligned}
    \right.
\end{equation}

\section{Figures, Tables and Algorithms}
Figures, Tables and Algorithms have to contain a Caption that describe their content, and have to be properly reffered in the text.

\subsection{Figures}
\label{subsec:figures}

For including pictures in your text you can use \texttt{TikZ} for high-quality hand-made figures,
or just include them as usual with the command
\begin{verbatim}
\includegraphics[options]{filename.xxx}
\end{verbatim}
Here xxx is the correct format, e.g. \verb|.png|, \verb|.jpg|, \verb|.eps|, \dots.

\begin{figure}[H]
    \centering
    \includegraphics[width=0.3\textwidth]{logo_polimi_scritta.eps}
    \caption{Caption of the Figure to appear in the List of Figures.}
    \label{fig:quadtree}
\end{figure}

Thanks to the \texttt{\textbackslash subfloat} command, a single figure, such as Figure~\ref{fig:quadtree},
can contain multiple sub-figures with their own caption and label, e.g. \color{black} Figure~\ref{fig:polimi_logo1} and Figure~\ref{fig:polimi_logo2}. 

\begin{figure}[H]
    \centering
    \subfloat[One PoliMi logo.\label{fig:polimi_logo1}]{
        \includegraphics[scale=0.5]{Images/logo_polimi_scritta.eps}
    }
    \quad
    \subfloat[Another one PoliMi logo.\label{fig:polimi_logo2}]{
        \includegraphics[scale=0.5]{Images/logo_polimi_scritta2.eps}
    }
    \caption[Shorter caption]{This is a very long caption you don't want to appear in the List of Figures.}
    \label{fig:quadtree2}
\end{figure}


\subsection{Tables}
\label{subsec:tables}

Within the environments \texttt{table} and  \texttt{tabular} you can create very fancy tables as the one shown in Table~\ref{table:example}.
\begin{table}[H]
    \caption*{\textbf{Title of Table (optional)}}
    \centering 
    \begin{tabular}{|p{3em} c c c |}
    \hline
    \rowcolor{bluepoli!40} % comment this line to remove the color
     & \textbf{column 1} & \textbf{column 2} & \textbf{column 3} \T\B \\
    \hline \hline
    \textbf{row 1} & 1 & 2 & 3 \T\B \\
    \textbf{row 2} & $\alpha$ & $\beta$ & $\gamma$ \T\B\\
    \textbf{row 3} & alpha & beta & gamma \B\\
    \hline
    \end{tabular}
    \\[10pt]
    \caption{Caption of the Table to appear in the List of Tables.}
    \label{table:example}
\end{table}

You can also consider to highlight selected columns or rows in order to make tables more readable.
Moreover, with the use of \texttt{table*} and the option \texttt{bp} it is possible to align them at the bottom of the page. One example is presented in Table~\ref{table:exampleC}. 

\begin{table}[H]
\centering 
    \begin{tabular}{|p{3em} | c | c | c | c | c | c|}
    \hline
%    \rowcolor{bluepoli!40}
     & \textbf{column1} & \textbf{column2} & \textbf{column3} & \textbf{column4} & \textbf{column5} & \textbf{column6} \T\B \\
    \hline \hline
    \textbf{row1} & 1 & 2 & 3 & 4 & 5 & 6 \T\B\\
    \textbf{row2} & a & b & c & d & e & f \T\B\\
    \textbf{row3} & $\alpha$ & $\beta$ & $\gamma$ & $\delta$ & $\phi$ & $\omega$ \T\B\\
    \textbf{row4} & alpha & beta & gamma & delta & phi & omega \B\\
    \hline
    \end{tabular}
    \\[10pt]
    \caption{Highlighting the columns}
    \label{table:exampleC}
\end{table}

\begin{table}[H]
\centering 
    \begin{tabular}{|p{3em} c c c c c c|}
    \hline
%    \rowcolor{bluepoli!40}
     & \textbf{column1} & \textbf{column2} & \textbf{column3} & \textbf{column4} & \textbf{column5} & \textbf{column6} \T\B \\
    \hline \hline
    \textbf{row1} & 1 & 2 & 3 & 4 & 5 & 6 \T\B\\
    \hline
    \textbf{row2} & a & b & c & d & e & f \T\B\\
    \hline
    \textbf{row3} & $\alpha$ & $\beta$ & $\gamma$ & $\delta$ & $\phi$ & $\omega$ \T\B\\
    \hline
    \textbf{row4} & alpha & beta & gamma & delta & phi & omega \B\\
    \hline
    \end{tabular}
    \\[10pt]
    \caption{Highlighting the rows}
    \label{table:exampleR}
\end{table}

\subsection{Algorithms}
\label{subsec:algorithms}

Pseudo-algorithms can be written in \LaTeX{} with the \texttt{algorithm} and \texttt{algorithmic} packages.
An example is shown in Algorithm~\ref{alg:var}.
\begin{algorithm}[H]
    \label{alg:example}
    \caption{Name of the Algorithm}
    \label{alg:var}
    \label{protocol1}
    \begin{algorithmic}[1]
    \STATE Initial instructions
    \FOR{$for-condition$}
    \STATE{Some instructions}
    \IF{$if-condition$}
    \STATE{Some other instructions}
    \ENDIF
    \ENDFOR
    \WHILE{$while-condition$}
    \STATE{Some further instructions}
    \ENDWHILE
    \STATE Final instructions
    \end{algorithmic}
\end{algorithm} 

\vspace{5mm}

\section{Theorems, propositions and lists}

\subsection{Theorems}
Theorems have to be formatted as:
\begin{theorem}
\label{a_theorem}
Write here your theorem. 
\end{theorem}
\textit{Proof.} If useful you can report here the proof.

\subsection{Propositions}
Propositions have to be formatted as:
\begin{proposition}
Write here your proposition.
\end{proposition}

\subsection{Lists}
How to  insert itemized lists:
\begin{itemize}
    \item first item;
    \item second item.
\end{itemize}
How to insert numbered lists:
\begin{enumerate}
    \item first item;
    \item second item.
\end{enumerate}

\section{Use of copyrighted material}

Each student is responsible for obtaining copyright permissions, if necessary, to include published material in the thesis.
This applies typically to third-party material published by someone else.

\section{Plagiarism}

You have to be sure to respect the rules on Copyright and avoid an involuntary plagiarism.
It is allowed to take other persons' ideas only if the author and his original work are clearly mentioned.
As stated in the Code of Ethics and Conduct, Politecnico di Milano \textit{promotes the integrity of research,
condemns manipulation and the infringement of intellectual property}, and gives opportunity to all those
who carry out research activities to have an adequate training on ethical conduct and integrity while doing research.
To be sure to respect the copyright rules, read the guides on Copyright legislation and citation styles available
at:
\begin{verbatim}
https://www.biblio.polimi.it/en/tools/courses-and-tutorials
\end{verbatim}
You can also attend the courses which are periodically organized on "Bibliographic citations and bibliography management".

\section{Bibliography and citations}
Your thesis must contain a suitable Bibliography which lists all the sources consulted on developing the work.
The list of references is placed at the end of the manuscript after the chapter containing the conclusions.
We suggest to use the BibTeX package and save the bibliographic references  in the file \verb|Thesis_bibliography.bib|.
This is indeed a database containing all the information about the references. To cite in your manuscript, use the \verb|\cite{}| command as follows:
\\
\textit{Here is how you cite bibliography entries: , or multiple ones at once:.}
\\
The bibliography and list of references are generated automatically by running BibTeX.

%-------------------------------------------------------------------------
%	BIBLIOGRAPHY
%-------------------------------------------------------------------------

\addtocontents{toc}{\vspace{2em}} % Add a gap in the Contents, for aesthetics
\bibliography{Thesis_bibliography} % The references information are stored in the file named "Thesis_bibliography.bib"

%-------------------------------------------------------------------------
%	APPENDICES
%-------------------------------------------------------------------------

\cleardoublepage
\addtocontents{toc}{\vspace{2em}} % Add a gap in the Contents, for aesthetics

% LIST OF FIGURES
\listoffigures

% LIST OF TABLES
\listoftables

\end{document}